\begin{claim}{D.9}{Proved}{Free hulls and reduction to injective free-monoid diagrams}
\label{D:finite-free-injective-reduction}
To decide whether every free-monoid diagram on
$A=\mathcal C_\delta$ passes its specified comparison, it is enough
to consider diagrams of finite-rank free monoids with injective
transition homomorphisms.

More precisely, for every diagram $M:A\to\Mon$ whose values are
finite-rank free monoids, there is a natural map $M\to M'$ such that
every $M'(a)$ is finite free, every transition of $M'$ is injective,
and the induced maps
\[
 B\!\int_A M\longrightarrow B\!\int_A M',
 \qquad
 \colim_A^{\Mon}M\longrightarrow\colim_A^{\Mon}M'
 \tag{D.9.1}
\]
are respectively a weak equivalence and an isomorphism. They commute
with the specified comparisons. Thus $q_M$ is a weak equivalence
if and only if $q_{M'}$ is.

The algebraic input is also valid independently of $A$: every finite
subset of a filtered colimit $S$ of free monoids lies in a finite-rank
free submonoid. Such a subset has a canonical least free submonoid
of $S$ containing it, of finite rank, and every automorphism of $S$
fixing the subset pointwise fixes that free hull pointwise.
\end{claim}
\begin{proof}
We use the free-hull and defect theorems for words
\citep[Section~2, Theorems~2.1--2.2]{DGCastiglioneFiciRestivo2019}
and the stability characterization of free submonoids
\citep[Theorem~1, p.~84]{DGCarpiDAlessandro2014}.
The latter says that a submonoid $L$ of a free monoid $F$ is free
if and only if, whenever $u\in F$ and
\[
 p,\ up,\ q,\ qu\in L,
 \tag{D.9.2}
\]
one has $u\in L$. We verify below exactly how these results apply
to the filtered colimit; no assertion that an arbitrary submonoid
of a free monoid is free is used.

\emph{Finite free charts in a filtered colimit.}
Write $S=\colim_i F_i$, with all $F_i$ free. It is cancellative:
an equality $xy=xz$ in $S$ is witnessed, together with representatives
of $x,y,z$, in one free stage; cancellation there makes $y=z$ in
$S$. Right cancellation is the same argument. It is also conical:
if $xy=1$, this equality is witnessed at a free stage, where both
factors are $1$.

Let $X\subset S$ be finite and lift its elements to a stage $F_i$.
Among all later stages and finite-rank free submonoids containing
the corresponding image of $X$, choose one, say $H\subset F_j$,
of minimum rank $r$. Such submonoids exist by the free-hull theorem,
and their ranks are nonnegative integers. If a later transition
were noninjective on $H$, the images of its $r$ basis letters would
satisfy a nontrivial word equation. An empty image or two coincident
images immediately reduces the number of generators to less than
$r$. Otherwise these are $r$ distinct nonempty words forming a
non-code; the defect theorem puts them in a free submonoid of rank
at most $r-1$. In either case the image of $X$ lies in a free
submonoid of smaller rank at that later stage, a contradiction.
Thus every later transition is injective on $H$, and $H\to S$ is
injective because equality in a filtered colimit is witnessed at
a later stage. Its image is the desired finite free chart. If
$X\subset\{1\}$, the trivial rank-zero chart already suffices.

\emph{An intrinsic finite free hull.}
Choose any finite free chart $H\subset S$ containing $X$ and form
the ordinary free hull $F_H(X)$ inside $H$. We show that it is the
least free submonoid of all of $S$ containing $X$, so in particular
it is independent of the chart.

If $L\subset S$ is any free submonoid, then $L\cap H$ is stable
inside $H$. Indeed, assume (D.9.2) holds in $L\cap H$, with
$u\in H$. In free $L$ the equality
\[
 q(up)=(qu)p
\]
has the usual Levi factorization. Either $qu=qt$ and $up=tp$
for some $t\in L$, in which case cancellation in $S$ gives
$t=u$, or $q=(qu)t$ for some $t\in L$, in which case cancellation
gives $ut=1$ and conicality gives $u=1$. Thus $u\in L\cap H$.
The stability criterion makes $L\cap H$ free. If $L$ contains
$X$, minimality of the ordinary free hull in $H$ gives
\[
 F_H(X)\subset L\cap H\subset L.
\]
Conversely $F_H(X)$ itself is a free submonoid of $S$ containing
$X$. This proves the asserted intrinsic minimality. Denote this
submonoid by $F_S(X)$. Its rank is at most $|X|$, by the free-hull
rank bound.

If an automorphism $g$ of $S$ fixes $X$ pointwise, minimality gives
$gF_S(X)=F_S(X)$. An automorphism of a free monoid permutes its
unique basis of indecomposable nonidentity elements. Every basis
letter of $F_S(X)$ occurs in some word of $X$: otherwise deleting
that letter from the basis gives a smaller free submonoid containing
$X$. The uniqueness of free factorization, applied to $g(x)=x$
for $x\in X$, therefore fixes each basis letter, hence the entire
hull, pointwise.

\emph{The replacement diagram.}
Use the concrete model of $A$ from
Claim~\ref{D:free-limit-criterion}: the objects are nonempty finite
closed substructures of the rich countable limit $U$, all closed
embeddings are arrows, $I$ is the directed inclusion poset, and
$G=\Aut(U)$. Form the smooth $G$-monoid
\[
 S=\colim_{B\in I}M(B).
\]
For each $B$, let $X_B\subset S$ be the image of the finite free
basis of $M(B)$ and set
\[
 M'(B)=F_S(X_B).
\]
This is finite free. A closed inclusion $B\subset D$ gives
$M'(B)\subset M'(D)$, since $M'(D)$ is free and contains the image
of $M(B)$. For any closed embedding $f:B\to D$, extend $f$ to
$g\in G$ and define $M'(f)$ by restricting $g$ on $S$. Equivariance
and minimality give
\[
 gM'(B)=F_S(gX_B)\subset M'(D).
\]
Two extensions differ, on the appropriate side, by an automorphism
fixing $B$ pointwise. Such an automorphism fixes $X_B$ and hence
$M'(B)$ pointwise by the hull argument. The restriction is therefore
independent of the extension. The same independence makes
composition strict. All these homomorphisms are injective, and the
canonical maps $M(B)\to S$ factor as a natural transformation
$M\to M'$.

The filtered union of the submonoids $M'(B)$ is exactly $S$:
it contains every image of $M(B)$ and is itself contained in $S$.
Consequently $M$ and $M'$ have the same smooth filtered-limit
monoid, with the natural map inducing its identity. The
product-compatible saturation comparison proved in
Claim~\ref{D:free-limit-criterion} now applies to the common diagram
$B\mapsto S^{G_B}$. Both $M$ and $M'$ map to this diagram, and
both maps induce weak equivalences on the homotopy colimits of
classifying spaces and isomorphisms on ordinary monoid colimits.
Their triangle commutes. Two-out-of-three gives (D.9.1), including
compatibility with the specified comparison maps.

\emph{Reduction of arbitrary ranks to finite ranks.}
Choose finitely many tagged basis letters $(B_i,s_i)$ of an
arbitrary free-monoid diagram $M$. At an object $D$, let $S_0(D)$
be the union of the basis supports of all words $M(f)(s_i)$,
where $f:B_i\to D$. This set is finite because each embedding set
$A(B_i,D)$ is finite. The free submonoids on the $S_0(D)$ form a
subdiagram: if a basis letter $t$ occurs in $M(f)(s_i)$, then
every letter of $M(g)(t)$ occurs in $M(gf)(s_i)$. Positivity of
word substitutions is essential here; there is no cancellation
that could remove that occurrence.

As the finite set of tagged seeds increases, these finite-rank
subdiagrams form a filtered union equal to $M$. Ordinary monoid
colimits and the Grothendieck-construction nerve commute with this
filtered union. For the latter, each simplex involves only finitely
many objects, arrows, and monoid elements; equivalently use the
simplicial replacement and the fact that filtered set colimits
commute with finite products. The same applies to the nerve of
the monoid colimit. Filtered colimits of weak equivalences of
simplicial sets are weak equivalences. Hence passage of the
comparison for every finite-rank subdiagram implies passage for
$M$. Combining this with (D.9.1) proves the stated sufficiency.
\end{proof}

\paragraph{What remains for $\mathcal C_\delta$.}
Claim~\ref{D:finite-free-injective-reduction} is a reduction of the
explicit question, not a resolution of it. An injective homomorphism
of free monoids need not induce an injective homomorphism of their
free-group completions. For example, the four length-two words
$aa,ab,ba,bb$ form a code, but if their abstract free generators are
$x,y,z,w$, then the nontrivial reduced word
$xz^{-1}wy^{-1}$ maps to $1$ in $F(a,b)$.

The injective replacement permits nonerasing word substitutions.
Their pairwise coincidence monoids are free: for
$f:P^*\to R^*$ and $g:Q^*\to R^*$ nonerasing, factor a pair
$(p,q)$ with $f(p)=g(q)$ at every common position that is both an
$f$-letter boundary and a $g$-letter boundary. The strictly increasing
boundary lists give a unique decomposition into minimal coincidence
pairs, proving freeness
\citep[Section~3, Lemma~3, revised version]{DGHolub2003}.
This controls a local overlap. A proof that these local ordered
refinements give all higher coherences in the comparison comma, or
an explicit injective free-monoid counterdiagram, is still needed.
